\chapter{Solusi latihan}
\label{eoceSolutions}





%_______________
\eocesolch{Pengantar data}



%_______________
\begin{multicols}{2}

% 1

\eocesol{(a)~Perlakuan: $10/43 = 0.23 \rightarrow 23\%$. \\
(b)~Kontrol: $2/46 = 0.04 \rightarrow 4\%$. 
(c)~Persentase pasien yang bebas nyeri 24 jam setelah menerima akupunktur
lebih tinggi dalam kelompok perlakuan.
(d)~Perbedaan persentase yang teramati antara kedua kelompok mungkin
terjadi karena kebetulan.}

% 3

\eocesol{(a)~``Apakah ada hubungan antara paparan polusi udara dan kelahiran prematur?''
(b)~143,196 kelahiran di California Selatan antara tahun 1989 dan 1993.
(c)~Pengukuran konsentrasi karbon monoksida, nitrogen dioksida, ozon, dan
partikulat kasar PM$_{10}$ (partikel dengan diameter aerodinamik hingga
10~$\mu m$, dengan konsentrasi dinyatakan dalam $\mu g/m^3$) yang dikumpulkan
di stasiun pemantau kualitas udara, serta lama kehamilan.
Semuanya merupakan variabel numerik kontinu. }

% 5

\eocesol{(a)~``Apakah memberi tahu anak-anak secara tegas agar tidak berbuat
curang memengaruhi kemungkinan mereka berbuat curang?''
(b)~160 anak berusia antara 5 dan 15 tahun.
(c)~Lima variabel: (1) usia (numerik, kontinu), (2) jenis kelamin (kategoris),
(3) apakah mereka anak tunggal (kategoris), (4) apakah mereka menerima petunjuk
tegas agar tidak berbuat curang (kategoris), dan (5) apakah mereka berbuat
curang (kategoris).}

% 7

\eocesol{Variabel penjelas: jenis akupunktur yang diterima, yaitu akupunktur
yang dirancang untuk menangani migrain atau akupunktur plasebo pada lokasi
yang bukan titik akupunktur.
Variabel respons: apakah pasien bebas nyeri setelah 24 jam.}

% 9

\eocesol{(a)~$50 \times 3 = 150$.
(b)~Empat variabel numerik kontinu: panjang sepal, lebar sepal, panjang petal, dan lebar petal.
(c)~Satu variabel kategoris, yaitu spesies, dengan tiga level: \emph{setosa}, \emph{versicolor}, dan \emph{virginica}.}

% 11

\eocesol{(a)~Status kepemilikan bandara (publik/privat),
    status penggunaan bandara (publik/privat),
    lintang,
    dan bujur.
(b)~Status kepemilikan bandara: kategoris, bukan ordinal.
    Status penggunaan bandara: kategoris, bukan ordinal.
    Lintang: numerik, kontinu.
    Bujur: numerik, kontinu.}

% 13

\eocesol{(a)~Populasi sasaran: kelahiran di California Selatan dalam cakupan
tempat dan waktu yang hendak diinferensikan; sampel: 143.196 kelahiran antara
tahun 1989 dan 1993 di California Selatan.
(b)~Jika kelahiran dalam sampel pada rentang waktu dan wilayah tersebut dapat
dianggap mewakili populasi sasaran, hasilnya dapat digeneralisasikan ke
kelahiran dalam cakupan itu. Namun, karena studi ini bersifat
observasional, temuannya tidak dapat digunakan untuk menetapkan hubungan 
kausal.}

% 15

\eocesol{(a)~Populasi: seluruh pasien asma berusia 18--69 tahun yang 
mengandalkan obat untuk pengobatan asma. Sampel: 600 pasien tersebut.
(b)~Jika pasien dalam sampel ini, yang kemungkinan tidak diambil secara acak, 
dapat dianggap mewakili seluruh pasien asma berusia 18--69 tahun yang 
mengandalkan obat untuk pengobatan asma, hasilnya dapat digeneralisasikan ke 
populasi yang didefinisikan di atas. Selain itu, karena studi ini merupakan 
eksperimen dengan penugasan acak, temuannya dapat digunakan untuk menetapkan 
hubungan kausal.}

% 17

\eocesol{(a)~Observasi.
(b)~Variabel.
(c)~Statistik sampel (rata-rata).
(d)~Parameter populasi (rata-rata).}

% 19

\eocesol{(a)~Studi observasional.
(b)~Gunakan pengambilan sampel berstrata untuk mengambil secara acak sejumlah 
tetap mahasiswa, misalnya 10 orang, dari setiap kelas praktikum sehingga 
ukuran sampel totalnya 40 mahasiswa.}

% 21

\eocesol{(a)~Positif, nonlinear, dan cukup kuat. Negara dengan persentase 
penduduk pengguna internet yang lebih tinggi juga cenderung memiliki harapan 
hidup saat lahir yang lebih tinggi, tetapi kenaikannya mulai melandai ketika 
mendekati 80 tahun.
(b)~Studi observasional.
(c)~Kemakmuran: negara yang penduduknya secara luas mampu membayar akses 
internet kemungkinan juga mampu membayar layanan kesehatan dasar. 
(Catatan: Jawaban dapat bervariasi.)}

% 23

\eocesol{(a)~Pengambilan sampel acak sederhana layak digunakan. Bahkan, jarang 
sekali metode ini tidak masuk akal untuk digunakan!
(b)~Pendapat mahasiswa mungkin berbeda menurut bidang studi, sehingga 
stratifikasi berdasarkan variabel ini masuk akal dan layak digunakan.
(c)~Mahasiswa yang seusia mungkin memiliki pendapat yang lebih serupa, padahal 
setiap klaster sebaiknya beragam sehubungan dengan hasil yang diteliti. Karena 
itu, pendekatan ini \textbf{tidak} baik. (Pertimbangan tambahan: jumlah orang 
dalam setiap klaster juga mungkin sangat berbeda sehingga ukuran sampel yang 
dihasilkan tidak terduga.)}

\end{multicols}
\begin{multicols}{2}

% 25

\eocesol{(a)~Kasusnya adalah 200 peserta dari kalangan laki-laki dan perempuan
yang diambil sebagai sampel acak.
(b)~Variabel responsnya adalah sikap terhadap sebuah oven gelombang mikro fiktif.
(c)~Variabel penjelasnya adalah sikap disposisional.
(d)~Ya, kasus-kasus tersebut diambil secara acak.
(e)~Ini adalah studi observasional karena tidak ada penugasan acak ke 
perlakuan.
(f)~Tidak, kita tidak dapat menetapkan hubungan kausal antara variabel 
penjelas dan variabel respons karena studi ini bersifat observasional.
(g)~Hasil hanya dapat digeneralisasikan ke populasi yang tercakup dalam
kerangka asal pemilihan 200 peserta. Soal tidak menyebutkan kerangka tersebut,
sehingga generalisasi ke populasi yang lebih luas tidak dapat
dipastikan.}

% 27

\eocesol{(a)~Sampel acak sederhana. Bias nonrespons dapat muncul: jika hanya 
orang-orang yang memiliki pendapat kuat tentang survei yang merespons, 
sampelnya mungkin tidak mewakili populasi.
(b)~Sampel kemudahan. Sampelnya mungkin tidak mewakili populasi karena hanya 
terdiri atas teman-temannya. Studi ini juga dapat mengalami bias nonrespons 
jika sebagian orang memilih untuk tidak mengembalikan survei.
(c)~Sampel kemudahan. Cara ini menghadapi masalah yang serupa dengan pembagian 
survei kepada teman-teman.
(d)~Pengambilan sampel multitahap. Bias bergantung pada peluang inklusi siswa,
pembobotan bila peluang itu berbeda, dan nonrespons. Kemiripan antarkelas
terutama memengaruhi presisi; kemiripan itu sendiri tidak menjamin bahwa
estimasi bebas bias.}

% 29

\eocesol{(a)~Kinerja ujian.
(b)~Kondisi pencahayaan: pencahayaan fluoresen dari atas, pencahayaan kuning dari atas, 
dan tanpa pencahayaan dari atas (hanya menggunakan lampu meja).
(c)~Variabel jenis kelamin: laki-laki, perempuan.}

% 31

\eocesol{(a)~Eksperimen.
(b)~Kondisi pencahayaan (pencahayaan fluoresen dari atas, pencahayaan kuning dari atas, 
dan tanpa pencahayaan dari atas yang hanya menggunakan lampu meja) serta tingkat 
kebisingan (tanpa kebisingan, kebisingan konstruksi, dan suara percakapan manusia).
(c)~Karena peneliti hendak memastikan bahwa laki-laki dan perempuan terwakili sama banyak 
dalam setiap perlakuan, variabel jenis kelamin menjadi variabel pemblokan.}

% 33

\eocesol{Diperlukan pengacakan dan penyamaran. Salah satu rancangan yang mungkin adalah: 
(1)~Siapkan dua cangkir yang identik untuk setiap peserta, satu berisi Coke biasa dan satu 
lagi berisi Diet Coke, dengan volume yang sama. Minta pihak yang tidak menyajikan minuman 
atau menilai respons untuk mengodekan identitas minuman secara acak sebagai A atau B pada 
setiap percobaan, lalu menyimpan kunci kode sampai respons selesai dicatat.
(2)~Sajikan kedua cangkir kepada setiap peserta satu per satu dalam urutan yang diacak, 
kemudian minta peserta mencatat nilai yang menunjukkan seberapa besar ia menyukai setiap 
minuman. Baik peserta maupun orang yang menyajikan cangkir tidak boleh mengetahui identitas 
minuman sampai seluruh penilaian selesai dicatat; dengan demikian eksperimen ini tersamar 
ganda. (Jawaban dapat bervariasi.)}

% 35

\eocesol{(a)~Studi observasional.
(b)~Anjing: Lucy. Kucing: Luna.
(c)~Oliver dan Lily.
(d)~Positif; ketika popularitas suatu nama untuk anjing meningkat,
popularitas nama tersebut untuk kucing juga cenderung meningkat.}

% 37

\eocesol{(a)~Eksperimen.
(b)~Perlakuan: 25 gram biji chia dua kali sehari; kontrol: plasebo.
(c)~Ya, jenis kelamin.
(d)~Ya, tersamar tunggal karena pasien tidak mengetahui perlakuan yang mereka terima.
(e)~Karena ini adalah eksperimen, kita dapat membuat pernyataan kausal. Namun, karena
sampelnya tidak acak, pernyataan kausal tersebut tidak dapat digeneralisasikan ke
populasi yang lebih luas.}

% 39

\eocesol{(a)~Orang yang tidak merespons mungkin akan memberikan jawaban berbeda
terhadap pertanyaan ini; misalnya, orang tua yang mengembalikan survei kemungkinan
memang tidak kesulitan meluangkan waktu bersama anak-anak mereka.
(b)~Kecil kemungkinan bahwa perempuan yang masih dapat dihubungi di alamat yang sama
tiga tahun kemudian merupakan sampel acak. Responden yang hilang ini mungkin lebih
banyak merupakan penyewa (bukan pemilik rumah), sehingga status sosial ekonominya
mungkin lebih rendah daripada responden yang berhasil dihubungi.
(c)~Studi ini tidak memiliki kelompok kontrol dan merupakan studi observasional;
selain itu, mungkin terdapat variabel perancu. Misalnya, orang-orang ini mungkin
berlari karena secara umum lebih sehat dan/atau juga melakukan olahraga lain.}

% 41

\eocesol{(a)~Eksperimen terkontrol acak.
(b)~Variabel penjelas: kelompok perlakuan (kategoris, dengan 3 level).
Variabel respons: kesejahteraan psikologis.
(c)~Tidak, karena para peserta adalah sukarelawan.
(d)~Ya, karena studi ini merupakan eksperimen.
(e)~Pernyataan tersebut seharusnya menggunakan kata ``bukti yang mendukung''
dan bukan ``membuktikan''.}

% 43

\eocesol{(a)~County, negara bagian, ras pengemudi, apakah kendaraan digeledah,
dan apakah pengemudi ditangkap.
(b)~Kelima variabel tersebut kategoris dan bukan ordinal.
(c)~Respons: apakah kendaraan digeledah atau tidak.
    Penjelas: ras pengemudi.}



%_______________
\end{multicols}



%_______________
\eocesolch{Merangkum data}



%_______________
\begin{multicols}{2}

% 1

\eocesol{(a)~Asosiasi positif: mamalia dengan masa kehamilan yang lebih lama cenderung memiliki 
rentang hidup yang lebih panjang.
(b)~Asosiasinya akan tetap positif.
(c)~Tidak, keduanya tidak saling bebas. Lihat bagian~(a).}

% 3

\eocesol{Grafik di bawah ini menunjukkan fase peningkatan.
Pada awalnya mungkin juga terdapat periode pertumbuhan eksponensial
sebelum ukuran cawan petri mulai memperlambat pertumbuhan. \\
\FigureFullPath[Sebuah grafik ditampilkan dengan sumbu horizontal "waktu" dan sumbu vertikal berlabel "jumlah sel bakteri". Sebuah kurva naik dengan curam di sisi kiri; ketika bergerak ke kanan, kenaikannya melambat hingga hampir berhenti ketika mencapai sisi kanan grafik.]{0.25}{ch_summarizing_data/figures/eoce/reproducing_bacteria/reproducing_bacteria_sketch}}

% 5

\eocesol{(a)~Rata-rata populasi, $\mu_{2007} = 52$; rata-rata sampel, $\bar{x}_{2008} = 58$.
(b)~Rata-rata populasi, $\mu_{2001} = 3.37$; rata-rata sampel, $\bar{x}_{2012} = 3.59$.}

% 7

\eocesol{10 karyawan mana saja yang rata-rata jumlah hari liburnya berada di antara jumlah hari 
libur minimum dan rata-rata jumlah hari libur seluruh pekerja di lokasi tersebut.}

% 9

\eocesol{(a)~Distribusi~2 memiliki rata-rata yang lebih tinggi karena $20 > 13$, dan simpangan 
baku yang lebih tinggi karena 20 lebih jauh dari data lainnya daripada 13.
(b)~Distribusi~1 memiliki rata-rata yang lebih tinggi karena $-20 > -40$, sedangkan 
Distribusi~2 memiliki simpangan baku yang lebih tinggi karena -40 lebih jauh dari data 
lainnya daripada -20.
(c)~Distribusi~2 memiliki rata-rata yang lebih tinggi karena semua nilai dalam distribusi 
ini lebih tinggi daripada nilai-nilai dalam Distribusi~1, tetapi kedua distribusi memiliki 
simpangan baku yang sama karena keduanya sama-sama bervariasi di sekitar rata-ratanya masing-masing.
(d)~Kedua distribusi memiliki rata-rata yang sama karena sama-sama berpusat pada 300, tetapi 
Distribusi~2 memiliki simpangan baku yang lebih tinggi karena pengamatannya lebih jauh dari 
rata-rata daripada pengamatan dalam Distribusi~1.}

% 11

\eocesol{(a)~Sekitar 30.
(b)~Karena distribusinya menceng kanan, rata-ratanya lebih tinggi daripada median.
(c)~Q1: antara 15 dan 20, Q3: antara 35 dan 40, IQR: sekitar 20.
(d)~Nilai yang dianggap sangat rendah atau sangat tinggi berada lebih dari 1.5$\times$IQR 
dari kuartil. Batas atas: Q3 + 1.5 $\times$ IQR = $37.5 + 1.5 \times 20 = 67.5$; 
batas bawah: Q1 - 1.5 $\times$ IQR = $17.5 - 1.5 \times 20 =  -12.5$. AQI terendah yang 
tercatat tidak lebih rendah daripada 5 dan AQI tertinggi yang tercatat tidak lebih tinggi 
daripada 65; keduanya berada di dalam batas tersebut. Oleh karena itu, tidak ada hari dalam 
sampel ini yang dapat dianggap memiliki AQI yang sangat rendah atau sangat tinggi.}

% 13

\eocesol{Histogram menunjukkan bahwa distribusinya bimodal, sesuatu yang tidak tampak pada 
diagram kotak. Diagram kotak memudahkan kita mengidentifikasi secara lebih tepat nilai-nilai 
pengamatan yang berada di luar garis kumis.}

% 15

\eocesol{(a)~Distribusi jumlah hewan peliharaan per rumah tangga kemungkinan menceng kanan 
karena terdapat batas alami pada 0 dan hanya sedikit orang yang memiliki banyak hewan 
peliharaan. Oleh karena itu, pusat distribusi paling tepat dideskripsikan dengan median dan 
variabilitasnya dengan IQR.
(b)~Distribusi jarak ke tempat kerja kemungkinan menceng kanan karena terdapat batas alami 
pada 0 dan hanya sedikit orang yang tinggal sangat jauh dari tempat kerja. Oleh karena itu, 
pusat distribusi paling tepat dideskripsikan dengan median dan variabilitasnya dengan IQR.
(c)~Distribusi tinggi badan laki-laki kemungkinan simetris. Oleh karena itu, pusat distribusi 
paling tepat dideskripsikan dengan rata-rata dan variabilitasnya dengan simpangan baku.}

% 17

\eocesol{(a)~Median merupakan ukuran yang jauh lebih baik untuk pendapatan tipikal 42 orang 
tersebut. Rata-rata jauh lebih tinggi daripada pendapatan 40 dari 42 orang itu. Hal ini 
terjadi karena rata-rata merupakan rerata aritmetika yang dipengaruhi oleh dua pengamatan 
ekstrem. Median tidak terlalu terpengaruh karena robust terhadap pencilan.
(b)~IQR merupakan ukuran yang jauh lebih baik untuk menggambarkan variabilitas pendapatan 
bagi hampir semua dari 42 orang tersebut. Simpangan baku sangat dipengaruhi oleh dua pendapatan tinggi, 
sedangkan IQR robust terhadap kedua pengamatan ekstrem itu.}

% 19

\eocesol{(a)~Distribusinya unimodal dan simetris, dengan rata-rata sekitar 25 menit dan 
simpangan baku sekitar 5 menit. Tampaknya tidak ada county dengan waktu perjalanan rata-rata 
yang sangat tinggi atau sangat rendah. Karena distribusinya sudah unimodal dan simetris, 
transformasi logaritma tidak diperlukan.
(b)~Jawaban dapat bervariasi. Terdapat kantong-kantong waktu perjalanan yang lebih lama di 
sekitar Washington, DC, bagian tenggara New York, Chicago, Minneapolis, Los Angeles, dan 
banyak kota besar lainnya. Terdapat pula wilayah luas dengan waktu perjalanan rata-rata yang 
lebih singkat dan bertumpang tindih dengan lahan pertanian di Midwest. Rumah banyak petani 
berdampingan dengan lahan pertaniannya sehingga perjalanan mereka ke tempat kerja singkat; 
hal ini dapat menjelaskan mengapa waktu perjalanan rata-rata di county-county tersebut 
relatif rendah.}

% 21

\eocesol{(a)~Diagram batang memperlihatkan urutan kategori dan frekuensi relatifnya.
(b)~Tidak ada ciri yang tampak pada diagram lingkaran tetapi tidak tampak pada diagram batang.
(c)~Pada umumnya kami lebih memilih diagram batang karena grafik ini juga memperlihatkan frekuensi relatif setiap kategori.}

% 23

\eocesol{Posisi vertikal tempat setiap kelompok ideologi terbagi ke dalam kategori Mendukung,
Tidak mendukung, dan Tidak yakin berbeda-beda. Hal ini menunjukkan bahwa peluang mendukung
DREAM Act berubah menurut ideologi politik, sehingga kedua variabel tersebut
mungkin saling dependen.}

\end{multicols}
\begin{multicols}{2}

% 25

\eocesol{(a)~(i) Salah. Alih-alih membandingkan jumlah, kita harus membandingkan persentase orang dalam setiap kelompok yang mengalami masalah kardiovaskular.
(ii)~Benar.
(iii)~Salah. Asosiasi tidak menyiratkan sebab-akibat. Kita tidak dapat menarik
kesimpulan kausal berdasarkan studi observasional. Perbedaannya dengan
bagian~(ii) tidak kentara.
(iv)~Benar. \\
(b)~Proporsi seluruh pasien yang mengalami masalah kardiovaskular: $\frac{7,979}{227,571} \approx 0.035$ \\
(c)~Jika kejadian masalah kardiovaskular dan perlakuan saling independen,
perkiraan jumlah pasien dengan masalah kardiovaskular dalam kelompok
rosiglitazon dapat dihitung dengan mengalikan jumlah pasien dalam kelompok itu
dengan tingkat kejadian masalah kardiovaskular secara keseluruhan dalam penelitian:
$67,593 * \frac{7,979}{227,571} \approx 2370$. \\
(d)~(i)~$H_0$: Perlakuan dan masalah kardiovaskular saling independen. Keduanya
tidak memiliki hubungan, dan perbedaan tingkat kejadian antara kelompok
rosiglitazon dan pioglitazon terjadi karena kebetulan.
$H_A$: Perlakuan dan masalah kardiovaskular tidak saling independen. Perbedaan
tingkat kejadian antara kelompok rosiglitazon dan pioglitazon tidak terjadi
karena kebetulan, dan rosiglitazon berasosiasi dengan peningkatan risiko
masalah kardiovaskular serius.
(ii)~Jumlah pasien dengan masalah kardiovaskular yang lebih tinggi daripada
perkiraan di bawah asumsi independensi akan mendukung hipotesis alternatif,
karena hal ini mengisyaratkan bahwa rosiglitazon meningkatkan risiko masalah
tersebut.
(iii)~Dalam penelitian sebenarnya, kita mengamati 2,593 kejadian kardiovaskular
dalam kelompok rosiglitazon. Pada 1,000 simulasi di bawah model independensi,
jumlah yang diamati dalam setiap simulasi selalu lebih kecil daripada 2,593.
Hal ini mengisyaratkan bahwa hasil sebenarnya tidak berasal dari model
independensi. Dengan kata lain, kedua variabel tampaknya tidak saling independen,
dan kita menolak model independensi serta memilih model alternatif. Hasil
penelitian memberikan bukti meyakinkan bahwa rosiglitazon berasosiasi dengan
peningkatan risiko masalah kardiovaskular.}


% 27

\eocesol{(a)~Menurunkan: nilai mahasiswa yang mengikuti ujian susulan lebih
kecil daripada rata-rata 24 nilai sebelumnya.
(b)~Hitung rata-rata tertimbang. Gunakan bobot 24 untuk rata-rata lama dan 1
untuk nilai baru: $(24\times 74 + 1\times64)/(24+1) = 73.6$.
%Ada cara lain untuk menyelesaikan latihan ini tanpa menggunakan rata-rata
%tertimbang.
(c)~Nilai baru berjarak lebih dari 1 simpangan baku dari rata-rata sebelumnya,
sehingga simpangan bakunya meningkat.}

% 29

\eocesol{Tidak. Kita memperkirakan distribusi ini menceng ke kanan karena dua
alasan: (1)~terdapat batas alami pada 0 (mustahil menonton televisi kurang dari
0 jam), dan (2)~simpangan baku distribusi sangat besar dibandingkan dengan
rata-ratanya.}

% 31

\eocesol{Distribusi usia pemenang kategori Aktris Terbaik menceng ke kanan dengan median
sekitar 30 tahun. Distribusi usia pemenang kategori Aktor Terbaik juga menceng ke kanan,
meskipun tidak sekuat itu, dengan median sekitar 40 tahun. Perbedaan letak
puncak kedua distribusi menunjukkan bahwa pemenang kategori Aktris Terbaik biasanya
lebih muda daripada pemenang kategori Aktor Terbaik. Usia pemenang kategori Aktris Terbaik lebih
bervariasi daripada usia pemenang kategori Aktor Terbaik. Kedua distribusi memiliki
kemungkinan pencilan pada ujung atasnya.}

% 33

\eocesol{\FigureFullPath[Ditampilkan diagram kotak untuk "Nilai" dengan kotak
membentang dari sekitar 72 hingga 82 dan median pada sekitar 78. Kumis
membentang ke bawah hingga 66 dan ke atas hingga 94. Satu titik ditampilkan di
bawah kumis bawah pada sekitar 57.]{0.25}{ch_summarizing_data/figures/eoce/stats_scores_box/stats_scores_boxplot}}



%_______________
\end{multicols}



%_______________
% Jawaban bernomor ganjil yang terhubung langsung dengan defining_probability.tex:
% hanya latihan 1, 3, 5, 7, 9, dan 11. Jawaban sumber adalah materi lapisan teks
% CC BY-SA; tidak ada solusi instruktur terbatas yang ditambahkan.
\eocesolch{Probabilitas}
\begin{multicols}{2}

% 1
\eocesol{(a)~Salah. Ini adalah percobaan independen.
(b)~Salah. Ada kartu bergambar berwarna merah.
(c)~Benar. Sebuah kartu tidak mungkin sekaligus kartu bergambar dan as.}

% 3
\eocesol{(a)~10 lemparan. Lebih sedikit lemparan berarti proporsi kepala dalam sampel lebih bervariasi,
sehingga peluang memperoleh setidaknya 60\% kepala lebih besar.
(b)~100 lemparan. Lebih banyak lemparan membuat proporsi kepala yang diamati sering lebih dekat
ke rata-rata, 0.50, dan karena itu juga di atas 0.40.
(c)~100 lemparan. Dengan lebih banyak lemparan, proporsi kepala yang diamati sering lebih dekat
ke rata-rata, 0.50.
(d)~10 lemparan. Lebih sedikit lemparan meningkatkan variasi fraksi lemparan yang menghasilkan kepala.}

% 5
\eocesol{(a)~$0.5^{10}$ = 0.00098.
(b)~$0.5^{10}$ = 0.00098.
(c)~$P$(setidaknya satu ekor) = $1 - P$(tidak ada ekor) = $1 - (0.5^{10}) \approx 1 - 0.001 = 0.999$.}

% 7
\eocesol{(a)~Tidak, ada pemilih yang sekaligus Independen dan pemilih mengambang. \\
(b)\\
\FigureFullPath[Diagram Venn untuk variabel “Independen” dan “Mengambang” ditampilkan; kedua lingkaran sebagian beririsan. Bagian lingkaran “Independen” yang tidak beririsan berlabel “24”. Bagian lingkaran “Mengambang” yang tidak beririsan berlabel “12”. Bagian irisan berlabel “11”.]{0.25}{ch_probability/figures/eoce/swing_voters/swing_voters.pdf} \\
(c)~Setiap pemilih Independen adalah pemilih mengambang atau bukan. Karena 35\% pemilih
adalah Independen dan 11\% adalah keduanya, 24\% sisanya pasti bukan pemilih mengambang.
(d)~0.47.
(e)~0.53.
(f)~P(Independen) $\times$ P(mengambang) = $0.35\times0.23 = 0.08$, yang
tidak sama dengan P(Independen dan mengambang) = 0.11, sehingga kedua peristiwa dependen.}

\end{multicols}
\begin{multicols}{2}

% 9
\eocesol{(a)~Jika kelas tidak dinilai berdasarkan kurva, keduanya independen. Jika dinilai
berdasarkan kurva, keduanya tidak independen maupun saling lepas--kecuali pengajar hanya akan
memberikan satu nilai A, situasi yang kita abaikan pada bagian~(b) dan~(c).
(b)~Keduanya mungkin tidak independen: jika Anda belajar bersama, kebiasaan belajar Anda
akan berkaitan, sehingga kinerja kuliah juga mungkin berkaitan.
(c)~Tidak. Lihat jawaban bagian~(a) ketika mata kuliah tidak dinilai berdasarkan kurva.
Secara umum, jika dua hal tidak berkaitan (independen), terjadinya salah satu tidak menghalangi
terjadinya yang lain.}

% 11
\eocesol{(a)~$0.16 + 0.09 = 0.25$.
(b)~$0.17 + 0.09 = 0.26$.
(c)~Dengan asumsi tingkat pendidikan suami dan istri independen:
$0.25 \times 0.26 = 0.065$. Perhatikan bahwa kita sebenarnya juga membuat asumsi kedua:
keputusan untuk menikah tidak berkaitan dengan tingkat pendidikan.
(d)~Asumsi independensi suami/istri mungkin tidak wajar karena orang sering menikah dengan
orang lain yang tingkat pendidikannya sebanding. Apakah asumsi kedua pada bagian~(c) wajar,
silakan Anda pertimbangkan.}

% Jawaban publik bernomor ganjil yang terhubung langsung dengan
% conditional_probability.tex: hanya latihan 13, 15, 17, 19, dan 21.
% Tidak ada jawaban instruktur terbatas yang ditambahkan. Latihan genap
% 14, 16, 18, 20, dan 22 dicatat sebagai kesenjangan penguasaan O001.
% Koreksi turunan: rujukan bagian pada jawaban 13(c) diperbaiki dari (a)
% menjadi (b-i), dan tata bahasa sumber pada jawaban 17(d) diperjelas.

% 13

\eocesol{(a)~Tidak, tetapi kita dapat menghitungnya jika A dan B saling independen.
(b-i)~0.21.
(b-ii)~0.79.
(b-iii)~0.3.
(c)~Tidak, karena 0.1 $\ne$ 0.21; nilai 0.21 dihitung berdasarkan independensi
pada bagian~(b-i).
(d)~0.143.}

% 15

\eocesol{(a)~Tidak, 0.18 responden termasuk dalam kombinasi ini.
(b)~$0.60 + 0.20 - 0.18 = 0.62$.
(c)~$0.18 / 0.20 = 0.9$.
(d)~$0.11 / 0.33 \approx 0.33$.
(e)~Tidak, karena jika saling independen, jawaban bagian~(c) dan~(d) akan sama.
(f)~$0.06 / 0.34 \approx 0.18$.}

% 17

\eocesol{(a)~Tidak. Ada 6~perempuan yang paling menyukai Five Guys Burgers.
(b)~$162 / 248 = 0.65$.
(c)~$181 / 252 = 0.72$.
(d)~Dengan asumsi independensi antara pemilihan pasangan kencan dan
preferensi hamburger--yang sepintas tampak wajar--diperoleh
$0.65 \times 0.72 = 0.468$.
(e)~$(252 + 6 - 1)/500 = 0.514$.}

% 19

\eocesol{(a) \\

\FigureFullPath[Diagram pohon dengan cabang utama “Dapat membuat diagram kotak?”
dan cabang sekunder “Lulus?”. Cabang utama “Dapat membuat diagram kotak” memiliki
dua kemungkinan, yaitu “Ya” dengan peluang 0.8 dan “Tidak” dengan peluang
0.2. Masing-masing cabang ini memiliki dua cabang sekunder. Cabang utama “Ya”
terbagi menjadi cabang “Ya” (untuk Lulus) dengan peluang bersyarat 0.86 dan
peluang akhir Ya-dan-Ya sebesar 0.688, serta cabang sekunder “Tidak” dengan
peluang bersyarat 0.14 dan peluang akhir Ya-dan-Tidak sebesar 0.112.
Cabang utama “Tidak” dari “Dapat membuat diagram kotak” memiliki cabang “Ya” dengan
peluang bersyarat 0.65 dan peluang akhir Tidak-dan-Ya sebesar 0.13, serta
cabang sekunder “Tidak” dengan peluang bersyarat 0.35 dan peluang akhir
Tidak-dan-Tidak sebesar 0.07.]{0.375}{ch_probability/figures/eoce/tree_drawing_box_plots/tree_drawing_box_plots}
(b)~0.84}

% 21

\eocesol{0.0714. Bahkan ketika pasien mendapat hasil tes positif untuk lupus,
peluang bahwa pasien tersebut benar-benar menderita lupus hanya 7,14\%.
House mungkin benar. \\

\FigureFullPath[Diagram pohon dengan cabang utama “Lupus” dan cabang sekunder
“Hasil” untuk tes lupus. Cabang utama “Lupus” memiliki dua kemungkinan, yaitu
“Ya” dengan peluang 0.02 dan “Tidak” dengan peluang 0.98. Masing-masing
cabang ini memiliki dua cabang sekunder. Cabang utama “Ya” terbagi menjadi cabang
“Ya” (untuk Hasil) dengan peluang bersyarat 0.98 dan peluang akhir
Ya-dan-Ya sebesar 0.0196, serta cabang sekunder “Tidak” dengan peluang
bersyarat 0.02 dan peluang akhir Ya-dan-Tidak sebesar 0.0004. Cabang utama
“Tidak” dari “Lupus” memiliki cabang sekunder “Ya” dengan peluang bersyarat
0.26 dan peluang akhir Tidak-dan-Ya sebesar 0.2548, serta cabang sekunder
“Tidak” dengan peluang bersyarat 0.74 dan peluang akhir Tidak-dan-Tidak
sebesar 0.7252.]{0.375}{ch_probability/figures/eoce/tree_lupus/tree_lupus.pdf}}

% Jawaban publik bernomor ganjil untuk Bab 3 Bagian 3.3 (latihan 23, 25, 27).
% Diambil hanya dari lapisan eoceSolutions publik yang dipatok; tidak ada jawaban terbatas yang ditambahkan.

% 23

\eocesol{(a)~0.3.
(b)~0.3.
(c)~0.3.
(d)~$0.3\times0.3=0.09$.
(e)~Ya, populasi yang menjadi sumber pengambilan sampel identik pada setiap pengambilan.}

% 25

\eocesol{(a)~$2 / 9 \approx 0.22$.
(b)~$3 / 9 \approx 0.33$.
(c)~$\frac{3}{10} \times \frac{2}{9} \approx 0.067$.
(d)~Tidak. Sebagai contoh, dalam latihan ini, diambilnya satu keping secara berarti
mengubah peluang hasil pengambilan berikutnya.}

% 27

\eocesol{$P(^1$legging, $^2$jins, $^3$jins$) = \frac{5}{24} \times \frac{7}{23} \times \frac{6}{22} = 0.0173$. 
Namun, mahasiswa yang mengenakan legging dapat terpilih pada urutan ke-2 atau ke-3; masing-masing 
urutan ini memiliki peluang yang sama, sehingga $3 \times 0.0173 = 0.0519$.}

% 29

\eocesol{(a)~13.
(b)~Tidak, 27 mahasiswa ini bukan sampel acak dari populasi mahasiswa 
universitas tersebut. Sebagai contoh, dapat dikatakan bahwa proporsi perokok di antara 
mahasiswa yang pergi ke pusat kebugaran pada pukul 9 pagi pada hari Sabtu lebih rendah daripada 
proporsi perokok di universitas secara keseluruhan.}

% 31

\eocesol{(a)~E(X) = 3.59. SD(X) = 9.64.
(b)~E(X) = -1.41. SD(X) = 9.64.
(c)~Tidak, laba bersih yang diharapkan bernilai negatif, sehingga secara rata-rata Anda diperkirakan akan kehilangan uang.}

% 33

\eocesol{Nilainya meningkat 5\%.}

% 35

\eocesol{E = -0.0526. SD = 0.9986.}

% 37

\eocesol{Jawaban perkiraan dapat diterima. \\
(a)~$(29+32)/144 = 0.42$.
(b)~$21/144 = 0.15$.
(c)~$(26+12+15)/144 = 0.37$.}

% Close untranslated Chapter 3 answer tail.
\end{multicols}

\eocesolch{Distribusi variabel acak}



%_______________
\begin{multicols}{2}

% 1

\eocesol{(a)~8.85\%.
(b)~6.94\%.
(c)~58.86\%.
(d)~4.56\%. \\
\FigureFullPath[Distribusi normal yang berpusat di 0 ditampilkan dengan ekor kiri kecil pada dan di bawah lokasi berlabel -1.35 diberi arsiran.]{0.23}{ch_distributions/figures/eoce/area_under_curve_1/zltNeg}
\FigureFullPath[Distribusi normal yang berpusat di 0 ditampilkan dengan ekor kanan kecil pada dan di atas lokasi berlabel 1.48 diberi arsiran.]{0.23}{ch_distributions/figures/eoce/area_under_curve_1/zgtPos}
\FigureFullPath[Distribusi normal yang berpusat di 0 ditampilkan dengan daerah tengah diberi arsiran. Ekor kiri yang besar sampai sedikit di bawah rata-rata dan ekor kanan yang kecil tetap tidak diarsir.]{0.23}{ch_distributions/figures/eoce/area_under_curve_1/zBet}
\FigureFullPath[Distribusi normal yang berpusat di nol ditampilkan dengan kedua ekor di bawah -2 dan di atas 2 diberi arsiran.]{0.23}{ch_distributions/figures/eoce/area_under_curve_1/zgtAbs}}

% 3

\eocesol{(a)~Verbal: $N(\mu = 151, \sigma = 7)$, Kuantitatif: $N(\mu = 153, \sigma = 7.67)$.
(b)~$Z_{VR} = 1.29$, $Z_{QR} = 0.52$. \\
\FigureFullPath[Ditampilkan sebuah distribusi normal dengan 2 garis vertikal yang diberi tanda khusus. Garis pertama terletak sedikit di atas rata-rata distribusi pada Z sama dengan 0.52 dan diberi label “QR” untuk Penalaran Kuantitatif. Garis kedua terletak lebih jauh di atas rata-rata pada Z sama dengan 1.29 dan diberi label “VR” untuk Penalaran Verbal.]{0.3}{ch_distributions/figures/eoce/GRE_intro/gre_intro.pdf} \\
(c)~Skornya berada 1.29 simpangan baku di atas rata-rata pada bagian
Penalaran Verbal dan 0.52 simpangan baku di atas rata-rata pada bagian
Penalaran Kuantitatif.
(d)~Ia berprestasi lebih baik pada bagian Penalaran Verbal karena skor-Z-nya
lebih tinggi pada bagian tersebut.
(e)~$Perc_{VR} = 0.9007 \approx 90\%$, $Perc_{QR} = 0.6990 \approx 70\%$.
(f)~$100\% - 90\% = 10\%$ memperoleh skor lebih tinggi daripadanya pada VR,
dan $100\% - 70\% = 30\%$ memperoleh skor lebih tinggi daripadanya pada QR.
(g)~Kita tidak dapat membandingkan skor mentah karena kedua bagian menggunakan
skala yang berbeda. Perbandingan skor persentil lebih tepat ketika membandingkan
kinerjanya dengan peserta lain.
(h)~Jawaban bagian (b)--(c) tidak berubah karena skor-Z dapat dihitung untuk
distribusi yang tidak normal. Namun, kita tidak dapat menjawab bagian~(d)--(f)
karena tanpa model normal kita tidak dapat menggunakan tabel peluang normal
untuk menghitung peluang dan persentil.}

% 5

\eocesol{(a)~$Z = 0.84$, yang bersesuaian dengan skor sekitar 159 pada QR.
(b)~$Z = -0.52$, yang bersesuaian dengan skor sekitar 147 pada VR.}

% 7

\eocesol{(a)~$Z = 1.2$, $P(Z > 1.2) = 0.1151$. \\
(b)~$Z= -1.28 \to 70.6\degree$F atau lebih rendah.}

% 9

\eocesol{(a)~$N(25, 2.78)$.
(b)~$Z = 1.08$, $P(Z > 1.08) = 0.1401$.
(c)~Jawabannya sangat berdekatan karena yang berubah hanyalah satuan.
(Jawabannya sedikit berbeda karena 28\degree C sama dengan 82.4\degree F,
bukan tepat 83\degree F.)
(d)~Karena $IQR = Q3 - Q1$, pertama-tama kita perlu mencari $Q3$ dan $Q1$,
kemudian menghitung selisihnya. Ingat bahwa $Q3$ adalah persentil ke-$75$
dan $Q1$ adalah persentil ke-$25$ suatu distribusi.
Dengan kuantil yang belum dibulatkan, Q1 = 23.1264\degree C,
Q3 = 26.8736\degree C, sehingga IQR = 3.7472\degree C,
atau sekitar 3.75\degree C.}

% 11

\eocesol{(a)~Tidak. Kartu-kartu tersebut tidak saling independen. Sebagai
contoh, jika kartu pertama adalah as keriting, kartu kedua tidak mungkin
merupakan as keriting. Selain itu, terdapat banyak kategori hasil yang
mungkin sehingga kategori-kategori tersebut perlu disederhanakan.
(b)~Tidak. Ada enam hasil yang sedang dipertimbangkan. Distribusi Bernoulli
hanya mengizinkan dua hasil atau kategori. Perhatikan bahwa lemparan dadu
dapat menjadi percobaan Bernoulli jika kita menyederhanakannya menjadi dua
hasil, misalnya keluar angka 6 dan tidak keluar angka 6. Namun, rincian
semacam itu perlu ditetapkan.}

% 13

\eocesol{(a)~$0.875^2\times 0.125 = 0.096$.
(b)~$\mu=8$, $\sigma=7.48$.}

% 15

\eocesol{Jika ${p}$ adalah peluang sukses, rata-rata variabel acak Bernoulli
$X$ diberikan oleh \\
$\mu = E[X] = P(X = 0) \times 0 + P(X = 1) \times 1$ \\
$= (1 - p) \times 0 + p\times 1 = 0 + p = p$}
% 17

\eocesol{(a)~Syarat-syarat binomial terpenuhi:
(1)~Percobaan yang saling independen: Dalam sampel acak, apakah seorang berusia
18--20 tahun telah mengonsumsi alkohol atau tidak, tidak bergantung pada
apakah orang lain telah mengonsumsinya atau tidak.
(2)~Jumlah percobaan yang tetap: $n = 10$.
(3)~Hanya ada dua hasil pada setiap percobaan: Mengonsumsi atau tidak
mengonsumsi alkohol.
(4)~Peluang sukses sama untuk setiap percobaan: $p = 0.697$.
(b)~0.203.
(c)~0.203.
(d)~0.167.
(e)~0.997.}

% 19

\eocesol{(a)~$\mu = 35$, $\sigma = 3.24$
(b)~$Z = \frac{45 - 35}{3.24} = 3.09$. Nilai 45 berjarak lebih dari 3
simpangan baku dari rata-rata sehingga kita dapat menganggapnya sebagai
pengamatan yang tidak biasa. Karena itu, ya, kita akan terkejut.
(c)~Dengan menggunakan pendekatan normal terhadap distribusi binomial,
0.0010. Dengan koreksi 0.5,
0.0017.}

% 21

\eocesol{(a)~$1-0.75^3 = 0.5781$.
(b)~0.1406.
(c)~0.4219.
(d)~$1-0.25^3=0.9844$.}

% 23

\eocesol{(a)~Distribusi geometrik: 0.109.
(b)~Binomial: 0.219.
(c)~Binomial: 0.137.
(d)~$1-0.875^6=0.551$.
(e)~Geometrik: 0.084.
(f)~Dengan menggunakan distribusi binomial dengan $n = 6$ dan $p=0.75$,
kita memperoleh $\mu=4.5$, $\sigma=1.06$, dan $Z = 2.36$. Karena nilai ini
tidak berada dalam jarak 2 simpangan baku, pengamatan tersebut dapat dianggap
tidak biasa.}

% 25

\eocesol{(a)~$\stackrel{Anna}{1/5}\times\stackrel{Ben}{1/4}\times\stackrel{Carl}{1/3}\times\stackrel{Damian}{1/2}\times\stackrel{Eddy}{1/1} = 1/5!=1/120$.
(b)~Karena peluang-peluang tersebut harus berjumlah 1, harus ada $5!=120$ kemungkinan
urutan.
(c)~$8!=\text{40,320}$.}
% 27

\eocesol{(a)~Geometrik, 0.0804.
(b)~Binomial, 0.0322.
(c)~Binomial negatif, 0.0193.}

% 29

\eocesol{(a)~Binomial negatif dengan $n=4$ dan $p=0.55$, dengan sukses didefinisikan sebagai siswa perempuan. Kasus binomial negatif tepat digunakan karena hasil percobaan terakhir ditetapkan, sedangkan urutan 3 percobaan pertama tidak diketahui.
(b)~0.1838.
(c)~${3 \choose 1} = 3$.
(d)~Dalam model binomial tidak ada pembatasan pada hasil percobaan terakhir. Dalam model binomial negatif, hasil percobaan terakhir ditetapkan. Oleh karena itu, kita memperhatikan banyaknya cara mengurutkan $k - 1$ sukses lainnya dalam $n - 1$ percobaan pertama.}

\end{multicols}
\begin{multicols}{2}

% 31

\eocesol{(a)~Poisson dengan $\lambda=75$.
(b)~$\mu=\lambda=75$, $\sigma=\sqrt{\lambda} = 8.66$.
(c)~$Z=-1.73$. Karena 60 berada dalam jarak 2 simpangan baku dari rata-rata, nilai ini secara umum tidak akan dianggap tidak biasa. Perhatikan bahwa kita sering menggunakan kaidah praktis ini bahkan ketika model normal tidak berlaku.
(d)~Dengan menggunakan distribusi Poisson dengan $\lambda = 75$: 0.0402.}

% 33

\eocesol{(a)~$\frac{\lambda^k \times e^{-\lambda}}{k!}
      = \frac{6.5^5 \times e^{-6.5}}{5!}
      = 0.1454$ \\
(b)~Peluangnya adalah
    $0.0015 + 0.0098 + 0.0318 = 0.0431$
    (0.0430 jika tidak ada galat pembulatan). \\
(c)~Rata-rata banyaknya orang per mobil adalah $11.7 / 6.5 = 1.8$,
    yang berarti orang-orang datang dalam kelompok-kelompok kecil.
    Artinya, jika satu orang datang, ada kemungkinan
    ia membawa satu atau lebih orang lain di dalam
    kendaraannya.
    Ini berarti setiap individu (orang) tidak saling independen,
    meskipun kedatangan mobil saling independen,
    sehingga salah satu asumsi inti distribusi Poisson
    tidak terpenuhi.
    Dengan demikian, banyaknya orang yang datang antara
    pukul 14.00 dan 15.00 tidak akan mengikuti distribusi Poisson.}

% Close untranslated Chapter 4 answer tail.
\end{multicols}

%_______________
\eocesolch{Dasar-dasar inferensi}



%_______________
\begin{multicols}{2}

% 1

\eocesol{(a)~Rata-rata. Setiap mahasiswa melaporkan nilai numerik, yaitu banyaknya jam.
(b)~Rata-rata. Setiap mahasiswa melaporkan suatu bilangan berupa persentase, dan persentase-persentase tersebut dapat dirata-ratakan.
(c)~Proporsi. Setiap mahasiswa menjawab Ya atau Tidak, sehingga variabel ini bersifat kategoris dan kita menggunakan proporsi.
(d)~Rata-rata. Setiap mahasiswa melaporkan suatu bilangan berupa persentase, seperti pada bagian~(b).
(e)~Proporsi. Setiap mahasiswa melaporkan apakah ia memperkirakan akan memperoleh pekerjaan atau tidak, sehingga variabel ini bersifat kategoris dan kita menggunakan proporsi.}

% 3

\eocesol{(a)~Sampel diambil dari seluruh cip komputer yang diproduksi
    di pabrik tersebut selama minggu produksi itu.
    Kita mungkin tergoda untuk menggeneralisasikan hasil ini
    ke seluruh minggu produksi, tetapi kita perlu berhati-hati
    karena tingkat cacat dapat berubah seiring waktu.
(b)~Proporsi cip komputer yang diproduksi
    di pabrik tersebut selama minggu produksi itu
    yang mengalami cacat.
(c)~Estimasi parameter dengan menggunakan data:
    $\hat{p} = \frac{27}{212} = 0.127$.
(d)~\emph{Galat baku} (atau $SE$).
(e)~Hitung $SE$ dengan menggunakan
    $\hat{p} = 0.127$ sebagai pengganti $p$:
    $SE
      \approx \sqrt{\frac{\hat{p}(1 - \hat{p})}{n}}
      = \sqrt{\frac{0.127(1 - 0.127)}{212}}
      = 0.023$.
(f)~Galat baku adalah
    simpangan baku dari $\hat{p}$.
    Nilai 0,10 berjarak sekitar satu galat baku
    dari nilai yang diamati, sehingga tidak
    merupakan simpangan yang sangat tidak biasa.
    (Sebagai kaidah praktis, nilai yang berjarak lebih dari sekitar
    2 galat baku biasanya dianggap tidak biasa.)
    Insinyur tersebut tidak perlu terkejut.
(g)~Galat baku yang dihitung ulang dengan $p = 0.1$:
    $SE = \sqrt{\frac{0.1(1 - 0.1)}{212}}
      = 0.021$.
    Nilai ini tidak jauh berbeda,
    sebagaimana lazimnya ketika galat baku
    dihitung menggunakan proporsi-proporsi yang relatif serupa
    (dan kadang-kadang bahkan ketika proporsinya cukup berbeda!).}

% 5

\eocesol{(a)~Distribusi sampling.
(b)~Jika proporsi populasi berada dalam rentang 5--30\%,
    syarat sukses--gagal akan terpenuhi dan
    distribusi sampling akan simetris.
(c)~Kita menggunakan rumus galat baku:
    $SE
      = \sqrt{\frac{p (1 - p)}{n}}
      = \sqrt{\frac{0.08 (1 - 0.08)}{800}}
      = 0.0096$.
(d)~Galat baku.
(e)~Distribusi cenderung lebih bervariasi
    ketika jumlah pengamatan per sampel lebih sedikit.}

% 7

\eocesol{Ingat bahwa rumus umumnya adalah $\text{estimasi titik} \pm z^{\star} \times SE$.
Pertama, tentukan ketiga nilai yang berbeda. Estimasi titiknya adalah 45\%, 
$z^{\star} = 1.96$ untuk tingkat kepercayaan 95\%, dan $SE = 1.2\%$. Kemudian, substitusikan 
nilai-nilai tersebut ke dalam rumus:
$ 45\% \pm 1.96 \times 1.2\% \quad\to\quad (42.6\%, 47.4\%) $
Kita yakin 95\% bahwa proporsi orang dewasa AS yang hidup dengan satu atau lebih 
penyakit kronis berada antara 42.6\% dan 47.4\%.}

% 9

\eocesol{(a)~Salah. Interval kepercayaan menyediakan rentang nilai yang masuk akal, dan 
kadang-kadang nilai sebenarnya tidak tercakup. Interval kepercayaan 95\% ``tidak mencakup'' nilai sebenarnya dalam sekitar 5\% 
kasus.
(b)~Benar. Perhatikan bahwa uraian tersebut berfokus pada nilai populasi yang sebenarnya.
(c)~Benar. Jika kita memeriksa interval kepercayaan 95\% yang dihitung dalam
Latihan~\ref{chronic_illness_intro}, kita dapat melihat bahwa interval ini tidak mencakup 50\%.
Hal ini 
berarti bahwa dalam uji hipotesis, kita akan menolak hipotesis nol bahwa 
proporsinya adalah~0.5.
(d)~Salah. Galat baku menggambarkan ketidakpastian dalam estimasi keseluruhan 
akibat fluktuasi alami karena keacakan, bukan ketidakpastian yang berkaitan 
dengan tanggapan tiap individu.}

% 11

\eocesol{(a)~Salah. Estimasi titik selalu berada dalam interval kepercayaan,
dan ini merupakan penerapan interval kepercayaan yang tidak masuk akal pada estimasi titik
(karena estimasi titik, berdasarkan konstruksinya, tercantum di dalam interval kepercayaan).
(b)~Benar.
(c)~Salah. Interval kepercayaan tersebut tidak berkaitan dengan rata-rata sampel.
(d)~Salah. Agar lebih yakin bahwa kita mencakup parameter, kita memerlukan interval yang lebih 
lebar. Bayangkan bahwa kita memerlukan jaring yang lebih besar agar lebih yakin dapat menangkap ikan di 
danau yang keruh.
(e)~Benar. Penjelasan opsional: Hal ini benar karena model normal digunakan untuk 
memodelkan rata-rata sampel. Batas galat adalah separuh lebar interval, 
dan rata-rata sampel adalah titik tengah interval.
(f)~Salah. Dalam penghitungan galat baku, kita membagi simpangan 
baku dengan akar kuadrat ukuran sampel. Untuk mengurangi SE (atau batas 
galat) menjadi separuh, kita perlu mengambil sampel sebanyak $2^2 = 4$ kali jumlah orang dalam 
sampel awal.}

% 13

\eocesol{(a)~Para pengunjung berasal dari sampel acak sederhana,
    sehingga independensi terpenuhi.
    Syarat sukses--gagal juga terpenuhi,
    dengan 64 dan $752 - 64 = 688$ keduanya di atas 10.
    Oleh karena itu, kita dapat menggunakan distribusi normal untuk
    memodelkan $\hat{p}$ dan membuat interval kepercayaan.
(b)~Proporsi sampelnya adalah $\hat{p} = \frac{64}{752} = 0.085$.
    Galat bakunya adalah
    {\footnotesize\begin{align*}
    SE
      &= \sqrt{\frac{p (1 - p)}{n}}
      \approx \sqrt{\frac{\hat{p} (1 - \hat{p})}{n}} \\
      &= \sqrt{\frac{0.085 (1 - 0.085)}{752}}
      = 0.010
    \end{align*}}%
(c)~Untuk interval kepercayaan 90\%,
    gunakan $z^{\star} = 1.6449$.
    Interval kepercayaannya adalah
    $0.085 \pm 1.6449 \times 0.010 \to
    (0.0683, 0.1017)$.
    Kita yakin 90\% bahwa 6.83\% hingga 10.17\%
    pengunjung pertama kali situs tersebut akan mendaftar menggunakan
    desain baru.}

% 15

\eocesol{(a)~$H_0: p = 0.5$
    (Siswa yang nilainya meningkat bukan merupakan mayoritas ataupun minoritas.)
  $H_A: p \neq 0.5$
    (Siswa yang nilainya meningkat merupakan mayoritas atau minoritas.) \\
(b)~$H_0: \mu = 15$
    (Rata-rata waktu kerja yang digunakan setiap karyawan untuk kegiatan
    nonpekerjaan selama March Madness adalah 15 menit.)
  $H_A: \mu \neq 15$
    (Rata-rata waktu kerja yang digunakan setiap karyawan untuk kegiatan
    nonpekerjaan selama March Madness berbeda dari 15 menit.)}

% 17

\eocesol{(1)~Hipotesis harus membahas
    proporsi populasi ($p$), bukan proporsi sampel.
(2)~Hipotesis nol harus memuat tanda sama dengan.
(3)~Hipotesis alternatif harus memuat tanda tidak sama dengan,
    dan
(4)~harus merujuk pada nilai nol, $p_0 = 0.6$,
    bukan proporsi sampel yang diamati.
Susunan hipotesis yang benar adalah:
$H_0: p = 0.6$ dan
$H_A: p \neq 0.6$.}

% 19

\eocesol{(a)~Klaim ini masuk akal karena seluruh interval
    berada di atas 50\%.
(b)~Nilai 70\% berada di luar interval,
    sehingga kita memiliki bukti yang meyakinkan bahwa dugaan peneliti
    tersebut keliru.
(c)~Interval kepercayaan 90\% akan lebih sempit daripada
    interval kepercayaan 95\%.
    Tanpa menghitung intervalnya sekalipun,
    kita dapat mengetahui bahwa 70\% tidak akan berada di dalam interval,
    dan kita juga akan menolak dugaan peneliti tersebut berdasarkan
    tingkat kepercayaan 90\%.}

% 21

\eocesol{(i)~Susun hipotesis. $H_0$: $p = 0.5$, $H_A$: $p \neq 0.5$.
  Kita akan menggunakan tingkat signifikansi $\alpha = 0.05$.
(ii)~Periksa syarat: sampel acak sederhana memberikan independensi,
  dan syarat sukses--gagal terpenuhi karena untuk setiap kelompok,
  $0.5 \times 1000 = 500$, yaitu sekurang-kurangnya~10.
(iii)~Selanjutnya, kita menghitung:
  $SE = \sqrt{0.5 (1 - 0.5) / 1000} = 0.016$.
  $Z = \frac{0.42 - 0.5}{0.016} = -5$,
  yang memiliki luas satu ekor sekitar 0.0000003,
  sehingga nilai-p adalah dua kali luas satu ekor tersebut, yaitu
  0.0000006.
(iv)~Tarik kesimpulan:
  Karena nilai-p lebih kecil daripada $\alpha = 0.05$,
  kita menolak hipotesis nol dan menyimpulkan bahwa
  proporsi orang dewasa AS yang meyakini kenaikan upah minimum
  akan membantu perekonomian tidak sama dengan 50\%.
  Karena nilai yang diamati lebih kecil daripada 50\% dan kita telah
  menolak hipotesis nol, kita dapat menyimpulkan
  bahwa keyakinan tersebut dianut oleh kurang dari 50\% orang dewasa AS.
  (Sebagai rujukan, survei tersebut juga meneliti dukungan terhadap
  perubahan upah minimum, yang merupakan pertanyaan berbeda dari
  apakah perubahan itu akan membantu perekonomian.)}

% 23

\eocesol{Jika nilai-p adalah 0.05, artinya statistik ujinya akan bernilai
$Z = -1.96$ atau $Z = 1.96$.
Kita akan menunjukkan penghitungan untuk $Z = 1.96$.
Galat baku:
$SE = \sqrt{0.3 (1 - 0.3) / 90} = 0.048$.
Terakhir, susun rumus statistik uji dan selesaikan
untuk memperoleh $\hat{p}$:
$1.96 = \frac{\hat{p} - 0.3}{0.048}
  \to \hat{p} = 0.394$
Sebagai alternatif, jika $Z = -1.96$ digunakan: $\hat{p} = 0.206$.}

\end{multicols}
\newpage
\begin{multicols}{2}

% 25

\eocesol{(a)~$H_0$: Antidepresan tidak memengaruhi gejala
    fibromialgia.
    $H_A$: Antidepresan memengaruhi gejala
    fibromialgia (baik membantu maupun memperburuk).
(b)~Menyimpulkan bahwa antidepresan membantu atau memperburuk
    gejala fibromialgia padahal sebenarnya tidak menimbulkan kedua efek tersebut.
(c)~Menyimpulkan bahwa antidepresan tidak memengaruhi
    gejala fibromialgia padahal sebenarnya berpengaruh.}
\end{multicols}

\newpage

\eocesolch{Inferensi untuk data kategoris}

\begin{multicols}{2}

% 1

\eocesol{(a)~Salah. Syarat sukses--gagal tidak terpenuhi.
(b)~Benar. Syarat sukses--gagal tidak terpenuhi. Dalam sebagian besar sampel, kita 
memperkirakan $\hat{p}$ mendekati 0.08, yaitu proporsi populasi sebenarnya. 
Meskipun $\hat{p}$ dapat jauh di atas 0.08, nilainya dibatasi dari bawah oleh 0, sehingga 
distribusinya diperkirakan menceng ke kanan. Membuat grafik distribusi sampling akan 
mengonfirmasi dugaan ini.
(c)~Salah. $SE_{\hat{p}} = 0.0243$, dan $\hat{p} = 0.12$ hanya berjarak 
$\frac{0.12 - 0.08}{0.0243} = 1.65$ galat baku dari rata-rata, sehingga tidak 
akan dianggap tidak lazim.
(d)~Benar. $\hat{p}=0.12$ berjarak 2.32 galat baku dari rata-rata, yang 
sering dianggap tidak lazim.
(e)~Salah. Galat baku berkurang dengan faktor $1/\sqrt{2}$.}

% 3

\eocesol{(a)~Benar. Lihat penalaran pada 6.1(b).
(b)~Benar. Dalam rumus galat baku, kita mengambil akar kuadrat ukuran sampel.
(c)~Benar. Syarat independensi dan syarat sukses--gagal terpenuhi.
(d)~Benar. Syarat independensi dan syarat sukses--gagal terpenuhi.}

% 5

\eocesol{(a)~Salah. Interval kepercayaan disusun untuk mengestimasi proporsi populasi, 
bukan proporsi sampel.
(b)~Benar. IK 95\%: $82\%\ \pm\ 2\%$.
(c)~Benar. Berdasarkan definisi tingkat kepercayaan.
(d)~Benar. Melipatgandakan ukuran sampel menjadi empat kali mengurangi galat baku dan batas galat dengan faktor 
$1/\sqrt{4}$.
(e)~Benar. Seluruh IK 95\% berada di atas 50\%.}

% 7

\eocesol{Dengan sampel acak, syarat independensi 
terpenuhi. Syarat sukses--gagal juga terpenuhi. 
$ME = z^{\star} \sqrt{ \frac{\hat{p} (1-\hat{p})} {n} } 
= 1.96 \sqrt{ \frac{0.56 \times  0.44}{600} }= 0.0397 \approx 4\%$}

% 9

\eocesol{(a)~Tidak. Sampel tersebut hanya merepresentasikan siswa yang mengikuti SAT, dan survei ini 
juga dilakukan secara daring.
(b)~(0.5289, 0.5711). Kita yakin 90\% bahwa 53\% hingga 57\% siswa kelas akhir 
sekolah menengah yang mengikuti SAT cukup yakin bahwa mereka akan mengikuti 
program belajar di luar negeri ketika berkuliah.
(c)~Sebanyak 90\% dari sampel acak semacam itu akan menghasilkan interval kepercayaan 90\% 
yang memuat proporsi sebenarnya.
(d)~Ya. Seluruh interval berada di atas 50\%.}

% 11

\eocesol{(a)~Kita ingin memeriksa apakah terdapat mayoritas (atau minoritas),
    sehingga kita menggunakan hipotesis berikut:
    \begin{align*}
    &H_0: p = 0.5
    &&H_A: p \neq 0.5
    \end{align*}
    Kita memiliki proporsi sampel $\hat{p} = 0.55$
    dan ukuran sampel $n = 617$ Independen. \\
    Karena ini adalah sampel acak, syarat independensi
    terpenuhi.
    Syarat sukses--gagal juga terpenuhi:
    $617 \times 0.5$ dan $617 \times (1 - 0.5)$
    keduanya sekurang-kurangnya 10 (kita menggunakan proporsi nol
    $p_0 = 0.5$ untuk pemeriksaan ini dalam uji hipotesis
    satu proporsi). \\
    Oleh karena itu, kita dapat memodelkan $\hat{p}$ menggunakan
    distribusi normal dengan galat baku sebesar
    \begin{align*}
    SE = \sqrt{\frac{p(1 - p)}{n}}
      = 0.02
    \end{align*}
    (Kita menggunakan proporsi nol $p_0 = 0.5$
    untuk menghitung galat baku dalam
    uji hipotesis satu proporsi.)
    Selanjutnya, kita menghitung statistik uji:
    \begin{align*}
    Z = \frac{0.55 - 0.5}{0.02} = 2.5
    \end{align*}
    Hasil ini memberikan luas satu ekor sebesar 0.0062,
    dan nilai-p sebesar $2 \times 0.0062 = 0.0124$. \\
    Karena nilai-p lebih kecil daripada 0.05,
    kita menolak hipotesis nol.
    Kita memiliki bukti kuat bahwa dukungan
    berbeda dari 0.5, dan karena data
    memberikan estimasi titik di atas 0.5,
    kita memiliki bukti kuat yang mendukung
    klaim komentator televisi tersebut. \\
(b)~Tidak.
    Secara umum, kita memperkirakan uji hipotesis
    dan interval kepercayaan akan selaras,
    sehingga kita memperkirakan interval kepercayaan
    menunjukkan rentang nilai masuk akal yang
    seluruhnya berada di atas 0.5.
    Namun, jika tingkat kepercayaan tidak selaras
    (misalnya tingkat kepercayaan 99\%
    dan tingkat signifikansi $\alpha = 0.05$),
    secara umum hal ini tidak lagi berlaku.}

% 13

\eocesol{(a)~$H_0: p = 0.5$. $H_A: p \neq 0.5$.
Syarat independensi (sampel acak) terpenuhi,
begitu pula syarat sukses--gagal (dengan menggunakan $p_0 = 0.5$,
kita memperkirakan 40 sukses dan 40 gagal).
$Z = 2.91$ $\to$ luas satu ekornya adalah 0.0018,
sehingga nilai-p adalah 0.0036.
Karena nilai-p $< 0.05$, kita menolak hipotesis nol.
Karena kita menolak $H_0$ dan estimasi titik menunjukkan bahwa orang-orang tersebut
lebih baik daripada sekadar menebak secara acak,
kita dapat menyimpulkan bahwa tingkat keberhasilan mengidentifikasi 
soda dengan benar pada orang-orang tersebut secara signifikan lebih tinggi daripada
tingkat keberhasilan yang diperoleh hanya dengan menebak secara acak.
(b)~Jika sebenarnya orang tidak dapat membedakan soda diet dari soda biasa 
dan mereka menebak secara acak, probabilitas memperoleh
sampel acak yang terdiri atas 
80 orang dengan 53 orang atau lebih mengidentifikasi soda secara benar
(atau 53 orang atau lebih mengidentifikasi soda secara keliru)
adalah 0.0036.}

% 15

\eocesol{Karena tersedia proporsi sampel ($\hat{p} = 0.55$),
kita menggunakannya dalam penghitungan ukuran sampel.
Batas galat untuk interval kepercayaan 90\% adalah
$1.6449 \times SE = 1.6449 \times \sqrt{\frac{p(1 - p)}{n}}$.
Kita ingin nilai ini lebih kecil daripada 0.01, dengan menggunakan
$\hat{p}$ sebagai pengganti $p$:
\begin{align*}
1.6449 \times \sqrt{\frac{0.55(1 - 0.55)}{n}} \leq 0.01 \\
1.6449^2 \frac{0.55(1 - 0.55)}{0.01^2} \leq n
\end{align*}
Dari hasil ini, kita memperoleh bahwa $n$ harus sekurang-kurangnya 6697.}

% 17

\eocesol{Studi ini bukan eksperimen acak, dan tidak jelas apakah orang-orang akan 
dipengaruhi oleh perilaku rekan mereka. Dengan kata lain, independensi mungkin tidak 
terpenuhi. Selain itu, hanya terdapat 5 intervensi dalam skenario 
provokatif, sehingga syarat sukses--gagal tidak terpenuhi. Bahkan jika kita mempertimbangkan 
uji hipotesis dengan menggabungkan kedua proporsi, syarat sukses--gagal 
tetap tidak akan terpenuhi. Karena satu syarat diragukan dan syarat 
lainnya tidak terpenuhi, selisih proporsi sampel tidak akan mengikuti 
distribusi yang mendekati normal.}

% 19

\eocesol{(a)~Salah. Seluruh interval kepercayaan berada di atas 0.
(b)~Benar.
(c)~Benar.
(d)~Benar.
(e)~Salah. Nilai-nilainya hanya dinegasikan dan diurutkan ulang: (-0.06,-0.02).}

% 21

\eocesol{(a)~Galat baku:
    \begin{align*}
    SE
      = \sqrt{\frac{0.79(1 - 0.79)}{347} +
          \frac{0.55(1 - 0.55)}{617}}
      = 0.03
    \end{align*}
    Dengan menggunakan $z^{\star} = 1.96$, kita memperoleh:
    \begin{align*}
    0.79 - 0.55 \pm 1.96 \times 0.03
      \to (0.181, 0.299)
    \end{align*}
    Kita yakin pada tingkat 95\% bahwa proporsi
    anggota Partai Demokrat yang mendukung rencana tersebut lebih tinggi 18.1 hingga
    29.9 poin persentase daripada proporsi
    pemilih independen yang mendukung rencana tersebut.
(b)~Benar.}

% 23

\eocesol{(a)~Lulusan perguruan tinggi: 23.7\%. Bukan lulusan perguruan tinggi: 33.7\%.
(b)~Misalkan $p_{CG}$ dan $p_{NCG}$ menyatakan proporsi lulusan perguruan tinggi 
dan bukan lulusan perguruan tinggi yang menjawab ``tidak tahu''. 
$H_0: p_{CG} = p_{NCG}$. $H_A: p_{CG} \ne p_{NCG}$. Syarat independensi terpenuhi 
(sampel acak), dan syarat sukses--gagal, 
yang diperiksa dengan menggunakan proporsi gabungan 
($\hat{p}_{\textit{pool}} = 235/827 = 0.284$), juga terpenuhi. $Z = -3.16$ $\to$ 
nilai-p $= 0.0016$. Karena nilai-p sangat kecil, kita menolak $H_0$. Data 
memberikan bukti kuat bahwa proporsi lulusan perguruan tinggi yang tidak 
memiliki pendapat tentang isu ini berbeda dari proporsi responden yang bukan lulusan perguruan 
tinggi. Data juga menunjukkan bahwa lebih sedikit lulusan perguruan tinggi yang mengatakan ``tidak 
tahu'' dibandingkan responden yang bukan lulusan perguruan tinggi (yakni, data menunjukkan arahnya setelah kita 
menolak $H_0$).}

% 25

\eocesol{(a)~Lulusan perguruan tinggi: 35.2\%. Bukan lulusan perguruan tinggi: 33.9\%.
(b)~Misalkan $p_{CG}$ dan $p_{NCG}$ menyatakan proporsi
lulusan perguruan tinggi 
dan bukan lulusan perguruan tinggi yang mendukung pengeboran lepas pantai.
$H_0: p_{CG} = p_{NCG}$. 
$H_A: p_{CG} \ne p_{NCG}$. Syarat independensi terpenuhi (sampel acak),
dan syarat sukses--gagal, yang diperiksa 
dengan menggunakan proporsi gabungan ($\hat{p}_{\textit{pool}} = 286/827 = 0.346$), juga 
terpenuhi. $Z = 0.37$ $\to$ nilai-p $=0.7113$. Karena nilai-p 
$> \alpha$ (0.05), kita gagal menolak $H_0$. Data tidak memberikan bukti kuat 
adanya perbedaan antara proporsi lulusan perguruan tinggi
dan bukan lulusan perguruan tinggi yang mendukung pengeboran lepas pantai di California.}

% 27

\eocesol{Subskrip $_C$ menyatakan kelompok kontrol. Subskrip $_T$ menyatakan pengemudi truk. 
$H_0: p_C = p_T$. $H_A: p _C \ne p_T$. Syarat independensi terpenuhi (sampel 
acak), demikian pula syarat sukses--gagal, yang 
diperiksa dengan menggunakan proporsi gabungan ($\hat{p}_{\textit{pool}} = 70/495 = 0.141$). 
$Z = -1.65$ $\to$ nilai-p $ = 0.0989$. Karena nilai-p lebih besar daripada $\alpha = 0.05$, kita gagal 
menolak $H_0$. Data tidak memberikan bukti kuat bahwa tingkat kurang 
tidur berbeda antara pekerja di luar sektor transportasi dan pengemudi truk.}

% 29

\eocesol{(a)~Ringkasan studi:
\begin{center}\scriptsize
\begin{tabular}{l l c c c}
        &           & \multicolumn{2}{c}{\textit{Kegagalan virologis}}   &       \\
\cline{3-4}
        &           & Ya       & Tidak        & Total \\
\cline{2-5}
\multirow{2}{*}{\textit{Pengobatan}}     & Nevirapine    & 26   & 94 & 120   \\
        & Lopinavir & 10            & 110   & 120   \\
\cline{2-5}
        & Total     & 36            & 204   & 240
\end{tabular}
\end{center}
(b)~$H_0: p_N = p_L$. Tidak terdapat perbedaan tingkat kegagalan virologis antara 
kelompok nevirapine dan lopinavir. $H_A: p_N \ne p_L$. Terdapat 
perbedaan tingkat kegagalan virologis antara kelompok nevirapine dan lopinavir.
(c)~Penempatan acak digunakan, sehingga pengamatan dalam setiap kelompok saling 
independen. Jika pasien dalam studi mewakili pasien dalam 
populasi umum (sesuatu yang mustahil diperiksa dari informasi yang diberikan), 
temuan tersebut juga dapat digeneralisasi kepada populasi dengan yakin. Syarat 
sukses--gagal, yang diperiksa dengan menggunakan proporsi gabungan 
($\hat{p}_{pool} = 36/240 = 0.15$), terpenuhi. $Z = 2.89$ $\to$ nilai-p 
$=0.0039$.
Karena nilai-p kecil, kita menolak $H_0$.
Terdapat bukti kuat adanya perbedaan tingkat kegagalan virologis antara
kelompok nevirapine dan lopinavir.
Pengobatan dan kegagalan virologis tampaknya tidak saling independen.}
\end{multicols}
